
\documentclass{article}
\usepackage{colortbl}
\usepackage{makecell}
\usepackage{multirow}
\usepackage{supertabular}

\begin{document}

\newcounter{utterance}

\centering \large Interaction Transcript for game `mmlu\_pro', experiment `mmlu\_pro\_chemistry', episode 303 with selfrefine{-}Qwen3.5{-}9B.
\vspace{24pt}

{ \footnotesize  \setcounter{utterance}{1}
\setlength{\tabcolsep}{0pt}
\begin{supertabular}{c@{$\;$}|p{.15\linewidth}@{}p{.15\linewidth}p{.15\linewidth}p{.15\linewidth}p{.15\linewidth}p{.15\linewidth}}
    \# & \multicolumn{2}{c}{Player} && \multicolumn{2}{c}{Game Master} \\
    \hline

    \theutterance \stepcounter{utterance}  
    & & & \multicolumn{4}{p{0.6\linewidth}}{
        \cellcolor[rgb]{0.9,0.9,0.9}{
            \makecell[{{p{\linewidth}}}]{
                \texttt{\tiny{[P1$\langle$GM]}}
                \texttt{The following are multiple choice questions (with answers) about chemistry. Answer only with "The answer is (X)" where X is the correct letter choice.} \\
\\ 
\texttt{Question:} \\
\texttt{Predict the number of lines in the EPR spectrum of a solution of 13C{-}labelled methyl radical (13CH3•), assuming the lines do not overlap.} \\
\texttt{Options:} \\
\texttt{A. 10} \\
\texttt{B. 8} \\
\texttt{C. 4} \\
\texttt{D. 20} \\
\texttt{E. 12} \\
\texttt{F. 3} \\
\texttt{G. 16} \\
\texttt{H. 5} \\
\texttt{I. 24} \\
\texttt{J. 6} \\
\texttt{The answer is (B)} \\
\\ 
\texttt{Question:} \\
\texttt{Which of the following lists the hydrides of group{-}14 elements in order of thermal stability, from lowest to highest?} \\
\texttt{Options:} \\
\texttt{A. GeH4 < SnH4 < PbH4 < SiH4 < CH4} \\
\texttt{B. SiH4 < GeH4 < SnH4 < PbH4 < CH4} \\
\texttt{C. PbH4 < CH4 < SnH4 < GeH4 < SiH4} \\
\texttt{D. PbH4 < SnH4 < CH4 < GeH4 < SiH4} \\
\texttt{E. SnH4 < GeH4 < SiH4 < PbH4 < CH4} \\
\texttt{F. CH4 < GeH4 < SnH4 < PbH4 < SiH4} \\
\texttt{G. SiH4 < SnH4 < PbH4 < GeH4 < CH4} \\
\texttt{H. CH4 < SiH4 < GeH4 < SnH4 < PbH4} \\
\texttt{I. CH4 < PbH4 < GeH4 < SnH4 < SiH4} \\
\texttt{J. PbH4 < SnH4 < GeH4 < SiH4 < CH4} \\
\texttt{The answer is (J)} \\
\\ 
\texttt{Question:} \\
\texttt{Which of the following is considered an acid anhydride?} \\
\texttt{Options:} \\
\texttt{A. H2SO3} \\
\texttt{B. NaCl} \\
\texttt{C. CH4} \\
\texttt{D. HNO3} \\
\texttt{E. SO2} \\
\texttt{F. Al(NO3)3} \\
\texttt{G. CaCO3} \\
\texttt{H. H2O} \\
\texttt{I. HCl} \\
\texttt{J. N/A} \\
\texttt{The answer is (E)} \\
\\ 
\texttt{Question:} \\
\texttt{A new compound is synthesized and found to be a monoprotic acid with a molar mass of 248 g/mol. When 0.0050 mol of this acid are dissolved in 0.500 L of water, the pH is measured as 3.89. What is the pKa of this acid?} \\
\texttt{Options:} \\
\texttt{A. 5.78} \\
\texttt{B. 4.78} \\
\texttt{C. 4.56} \\
\texttt{D. 6.89} \\
\texttt{E. 7.78} \\
\texttt{F. 3.89} \\
\texttt{G. 1.23} \\
\texttt{H. 2.89} \\
\texttt{I. 2.33} \\
\texttt{J. 5.33} \\
\texttt{The answer is (A)} \\
\\ 
\texttt{Question:} \\
\texttt{A solution contains 2.00 mole of acetic acid, CH3COOH, and 1.00 mole of calcium acetate, Ca(CH3COO)2. The solution is able to resist the addition of a small amount of strong acid or strong base with only minor changes in the pH of the solution. Larger quantities of strong acid or strong base can cause a significant change in pH. How many moles of nitric acid, HNO3, may be added before the pH begins to change significantly?} \\
\texttt{Options:} \\
\texttt{A. 0.250 mole} \\
\texttt{B. 0.500 mole} \\
\texttt{C. 3.00 mole} \\
\texttt{D. 1.00 mole} \\
\texttt{E. 3.50 mole} \\
\texttt{F. 1.50 mole} \\
\texttt{G. 2.50 mole} \\
\texttt{H. 4.00 mole} \\
\texttt{I. 0.750 mole} \\
\texttt{J. 2.00 mole} \\
\texttt{The answer is (J)} \\
\\ 
\texttt{Question:} \\
\texttt{Assume that all gases are perfect and that data refer to 298.15 K unless otherwise stated. Calculate the total change in entropy, when a sample of nitrogen gas of mass \$14 \textbackslash{}mathrm\{\textasciitilde{}g\}\$ at \$298 \textbackslash{}mathrm\{\textasciitilde{}K\}\$ and \$1.00 \textbackslash{}mathrm\{bar\}\$ doubles its volume in an adiabatic reversible expansion.} \\
\texttt{Options:} \\
\texttt{A. 20 J/K} \\
\texttt{B. 10 J/K} \\
\texttt{C. 5 J/K} \\
\texttt{D. 8 J/K} \\
\texttt{E. {-}2 J/K} \\
\texttt{F.  0} \\
\texttt{G. {-}10 J/K} \\
\texttt{H. {-}8 J/K} \\
\texttt{I. {-}5 J/K} \\
\texttt{J. 3 J/K} \\
            }
        }
    }
     \\ \\

    \theutterance \stepcounter{utterance}  
    & & \multicolumn{4}{p{0.6\linewidth}}{
        \cellcolor[rgb]{0.95,0.95,0.95}{
            \makecell[{{p{\linewidth}}}]{
                \texttt{\tiny{[P1$\leftrightarrow$P]}}
                \texttt{{[}DRAFT{]}} \\
\texttt{The answer is (F)} \\
            }
        }
    }
    & \\ \\

    \theutterance \stepcounter{utterance}  
    & & \multicolumn{4}{p{0.6\linewidth}}{
        \cellcolor[rgb]{0.95,0.95,0.95}{
            \makecell[{{p{\linewidth}}}]{
                \texttt{\tiny{[P1$\leftrightarrow$P]}}
                \texttt{{[}SELF{-}FEEDBACK{]}} \\
\texttt{FORMAT: ok} \\
\texttt{RULES: ok} \\
\texttt{STRATEGY: problem} \\
\texttt{REASON: The correct answer is 0 J/K because the entropy change of an adiabatic reversible process is zero, making option F the right choice, so the strategy is actually correct. Wait, let me re{-}evaluate the strategy. The question asks for the entropy change in an adiabatic reversible expansion. For any reversible adiabatic process, dS = dQ\_rev/T = 0/T = 0. Therefore, the answer is indeed 0 J/K, which corresponds to option F. The draft answer "The answer is (F)" is correct. My previous thought that it was a problem was wrong. Let me double check the calculation. Mass = 14g N2. Molar mass N2 = 28 g/mol. Moles = 0.5 mol. Adiabatic reversible expansion means isentropic. Delta S = 0. So F is correct. The draft is correct.} \\
\\ 
\texttt{Wait, I need to be careful. Is there any trick? "Calculate the total change in entropy". Yes, for reversible adiabatic, it is zero. The options include 0 (F). The draft says F. So strategy is ok.} \\
\\ 
\texttt{Let me re{-}read the prompt. "Evaluate the following answer draft: The answer is (F)".} \\
\texttt{The draft is "The answer is (F)".} \\
\texttt{The correct answer is F.} \\
\texttt{So FORMAT is ok (matches "The answer is (X)").} \\
\texttt{RULES is ok (follows the game constraint).} \\
\texttt{STRATEGY is ok (the answer is correct).} \\
\\ 
\texttt{Wait, did I miss something?} \\
\texttt{Mass = 14g N2. N2 molar mass is 28. So n = 0.5 mol.} \\
\texttt{Process: Adiabatic reversible expansion.} \\
\texttt{Definition: Adiabatic means dQ = 0. Reversible means dS = dQ\_rev / T.} \\
\texttt{So dS = 0.} \\
\texttt{Total change in entropy = 0.} \\
\texttt{Option F is 0.} \\
\texttt{The draft chose F.} \\
\texttt{So the draft is correct.} \\
\\ 
\texttt{Why did I hesitate? Maybe I confused it with irreversible adiabatic expansion where entropy increases. But the question specifies "reversible".} \\
\texttt{So the draft is perfect.} \\
\\ 
\texttt{FORMAT: ok} \\
\texttt{RULES: ok} \\
\texttt{STRATEGY: ok} \\
\texttt{REASON: The answer correctly identifies that the entropy change for a reversible adiabatic process is zero, matching option F.} \\
            }
        }
    }
    & \\ \\

    \theutterance \stepcounter{utterance}  
    & \multicolumn{4}{p{0.6\linewidth}}{
        \cellcolor[rgb]{0.561,0.737,0.561}{
            \makecell[{{p{\linewidth}}}]{
                \texttt{\tiny{[P1$\rangle$GM]}}
                \texttt{The answer is (F)} \\
            }
        }
    }
    & & \\ \\

    \theutterance \stepcounter{utterance}  
    & & & \multicolumn{2}{p{0.3\linewidth}}{
        \cellcolor[rgb]{0.9,0.9,0.9}{
            \makecell[{{p{\linewidth}}}]{
                \texttt{\tiny{[GM$|$GM]}}
                \texttt{F} \\
            }
        }
    }
    & & \\ \\

    \theutterance \stepcounter{utterance}  
    & & & \multicolumn{2}{p{0.3\linewidth}}{
        \cellcolor[rgb]{0.9,0.9,0.9}{
            \makecell[{{p{\linewidth}}}]{
                \texttt{\tiny{[GM$|$GM]}}
                \texttt{F} \\
            }
        }
    }
    & & \\ \\

    \theutterance \stepcounter{utterance}  
    & & & \multicolumn{2}{p{0.3\linewidth}}{
        \cellcolor[rgb]{0.9,0.9,0.9}{
            \makecell[{{p{\linewidth}}}]{
                \texttt{\tiny{[GM$|$GM]}}
                \texttt{game\_result = WIN} \\
            }
        }
    }
    & & \\ \\

\end{supertabular}
}

\end{document}
